% =============================================================================
% main.tex — ESG Risk and Regulatory Uncertainty: Evidence from U.S.
%            Climate-Disclosure Shocks
% Byeong-Hak Choe (SUNY Geneseo). Single-authored draft.
%
% Tables (\input) and figures (\includegraphics) are generated by the R
% pipeline into scripts/R/_outputs/ (gitignored). Run the pipeline first:
%   Rscript scripts/R/00_run_all.R
% then compile:
%   cd paper && pdflatex -interaction=nonstopmode main.tex
%
% Citations resolve against ../Bibliography_base.bib (plainnat). Compile:
%   pdflatex main; bibtex main; pdflatex main; pdflatex main
% Remaining \TODO markers are CONTENT placeholders (acknowledgments, an
% implication sentence, a magnitude benchmark), not citations.
% =============================================================================
\documentclass[12pt]{article}

\usepackage[margin=1in]{geometry}
\usepackage{amsmath, amssymb}
\usepackage{graphicx}
\usepackage{booktabs}
\usepackage{threeparttable}
\usepackage{caption}
\usepackage[hidelinks]{hyperref}
\usepackage{setspace}
\usepackage{xcolor}
\usepackage[authoryear,round]{natbib}

\graphicspath{{../scripts/R/_outputs/}}
\newcommand{\TODO}[1]{\textcolor{red}{[TODO: #1]}}
\captionsetup{font=small, labelfont=bf}
\onehalfspacing

\title{\textbf{ESG Risk and Regulatory Uncertainty:\\
Evidence from U.S.\ Climate-Disclosure Shocks}}
\author{Byeong-Hak Choe\thanks{SUNY Geneseo. Email: \texttt{choe@geneseo.edu}.
This research uses ESG Risk Ratings from Morningstar Sustainalytics and financial
data from Yahoo!\ Finance; all tables and figures are generated by a reproducible
R pipeline (\texttt{scripts/R/}). I thank colleagues and seminar participants for
helpful comments. All remaining errors are my own.}}
\date{This draft: \today}

\begin{document}
\maketitle

\begin{abstract}
\noindent How do firms' measured environmental, social, and governance (ESG) risk profiles
respond when climate-disclosure policy becomes both more salient and more uncertain? I study
U.S.\ public companies around the U.S.\ Securities and Exchange Commission's 2024
climate-disclosure rule, a major but legally contested mandate that sharply raised uncertainty
about future climate-risk reporting. Using firm-level Sustainalytics ESG risk ratings from
March 2024 and March 2025, I estimate within-firm changes in measured E, S, and G risk and
instrument those changes with firms' exposure to the contemporaneous Sustainalytics
rating-system overhaul, isolating its common industry-level component as a shift-share
instrument. I document two facts. First, measured \emph{governance} risk fell for the large
majority of firms over this window (88.9\% of firms; mean change $-1.19$ points), while
environmental risk drifted modestly upward and controversy scores were largely unchanged.
Second, instrumenting the change in measured S and G risk (where the instrument is strong for
governance and moderate for social risk, with first-stage $F = 15.0$ and $9.5$, respectively), I
find no statistically detectable pass-through to controversy, leverage, investment, debt issuance,
or buybacks within the first year; weak-instrument-robust Anderson--Rubin intervals confirm the
nulls. The average null nonetheless masks heterogeneity: changes in environmental risk are
associated with real financing and investment behavior precisely for the firms the rule binds on
(large filers and climate-sensitive sectors), with several interactions surviving a multiple-testing
correction. The
immediate footprint of the disclosure-rule episode on measured ESG risk was thus large but
concentrated in the rating methodology, with real-behavior responses confined to the most exposed
firms. A wholesale movement in measured ESG risk should therefore not be read as evidence of
changed corporate behavior, especially around a contested disclosure mandate.
\end{abstract}

% =============================================================================
\section{Introduction}
% =============================================================================
In March 2024 the U.S.\ Securities and Exchange Commission adopted its long-awaited
climate-disclosure rule, requiring public companies to report climate-related risks and, for
larger filers, certain greenhouse-gas emissions. Within weeks the rule was stayed amid
litigation, leaving firms and investors facing an unusual combination of high salience and deep
uncertainty about what climate-risk reporting would ultimately require. This episode is a
natural setting in which to ask a simple question: when climate-disclosure policy becomes more
prominent yet more uncertain, do firms' \emph{measured} ESG risk profiles actually move, and
does any such movement pass through to corporate and capital-market outcomes?

Answering this question is difficult because ESG ratings are observed in the cross-section, where
they correlate with size, sector, profitability, and a host of unobserved firm attributes.
Cross-sectional ESG correlations therefore tell us little about whether a \emph{change} in
measured ESG risk has consequences. I make progress by exploiting two features of this period.
First, I observe the same firms before and after the rule, in March 2024 and March 2025, so I
can difference out fixed firm attributes and study within-firm changes. Second, the rating
provider Sustainalytics implemented a methodology overhaul over this same window that mechanically
revised measured E, S, and G risk in ways that varied across industries. I use firms' exposure to
that rating-system change as an instrument for the change in their measured pillar risk, isolating
its common industry-level component as a shift-share (leave-one-out industry-mean) instrument.

I find, first, that the within-firm footprint of this period on measured ESG risk was large but
uneven across pillars. Governance risk fell for 88.9\% of firms (mean change $-1.19$ points on
the Sustainalytics scale, where higher denotes more risk), environmental risk rose modestly
(55.6\% of firms; mean $+0.38$), and controversy scores were essentially unchanged for four firms
in five. Second, where the instrument is informative---Social ($F = 9.5$) and Governance
($F = 15.0$)---two-stage least squares yields no statistically detectable effect of the change in
measured pillar risk on controversy, leverage, asset growth, capital expenditure, net debt
issuance, or share buybacks within the first year. Because the first stages are only moderately
strong, I report Anderson--Rubin confidence intervals that are robust to weak identification;
they are wider than the conventional intervals but continue to include zero. Environmental risk
cannot be instrumented strongly in these data ($F = 5.2$) and is reported as an ordinary-least-
squares association only.

The paper makes two contributions. First, it moves beyond cross-sectional ESG correlations by
identifying \emph{within-firm} changes in measured ESG risk and pairing them with a policy-driven
rating shock, which lets it separate methodology-driven movements in measured risk from changes in
firm fundamentals. Second, it shows that the immediate measured-ESG footprint of the 2024
disclosure-rule episode was concentrated in governance and, plausibly, in the rating methodology
itself, with no detectable short-run transmission to real corporate behavior on average and a
response confined to the most rule-exposed firms.

The analysis connects to three literatures. The first studies the \emph{measurement} of ESG: ratings
diverge sharply across providers \citep{BergKolbelRigobon2022_aggregate_confusion} and are
retroactively revised \citep{BergFabisikSautner2021_rewriting_history}, which is why I study
within-firm changes in a single provider rather than cross-sectional levels. Most closely,
\citet{RzeznikWeissHanleyPelizzon2026_investor_reliance} exploit Sustainalytics' 2018 methodology
overhaul as an exogenous shock to measured ESG risk and trace its effect on prices through investor
reliance; I use a later vintage of the same provider's methodology change as an instrument and study
real corporate behavior rather than the pricing channel. The second is the literature on
disclosure regulation \citep{ChristensenHailLeuz2021_mandatory_csr}, including market reactions to
the SEC climate-disclosure proposal \citep{Kim2024_sec_climate_reaction}. The third asks whether
ESG and climate risk are priced \citep{KruegerSautnerStarks2020_climate_institutional,
BoltonKacperczyk2021_carbon_risk, PastorStambaughTaylor2021_sustainable_equilibrium}, the
market-outcome margin I flag as the first-order open question.

The remainder of the paper proceeds as follows. Section~\ref{sec:data} describes the policy
setting and the data. Section~\ref{sec:strategy} lays out the first-difference design, the
shift-share instrument, and its identifying assumptions. Section~\ref{sec:results} presents the
results and weak-instrument-robust inference. Section~\ref{sec:conclusion} concludes.

% =============================================================================
\section{Institutional Background and Data}\label{sec:data}
% =============================================================================
\paragraph{The 2024 climate-disclosure rule.} The SEC's March 2024 rule was the first
federal mandate requiring climate-risk disclosure by U.S.\ public companies. Its near-immediate
judicial stay meant that, throughout my sample window, firms operated under a regime of elevated
attention to climate risk combined with genuine uncertainty about the eventual reporting
standard.\footnote{The Commission adopted the final rule on March 6, 2024 (Release No.\ 33-11275)
and stayed it on April 4, 2024 pending judicial review (Release Nos.\ 33-11280 and 34-99908).
\citet{Kim2024_sec_climate_reaction}
documents a negative market reaction to the earlier March 2022 proposal.}

\paragraph{The Sustainalytics rating-system change.} Over the same period, Sustainalytics
revised its ESG Risk Ratings methodology, changing how Material ESG Issues are measured and
weighted at the subindustry level. Because these revisions applied unevenly across industries,
they generated variation in the measured change in E, S, and G risk that is plausibly orthogonal
to a given firm's idiosyncratic fundamentals---variation I exploit below.

\paragraph{Data and sample.} I combine firm-level Sustainalytics ESG risk ratings for U.S.\
companies from the March 2024 and March 2025 snapshots with annual financial statements
(income statement, balance sheet, and cash-flow statement) used to construct controls and real-behavior
outcomes. The estimation sample is the balanced panel of 625 firms present in both ESG snapshots;
613 of these also have usable financial data. All financial controls are measured at the most
recent fiscal year ending before the March 2024 snapshot, and real-behavior outcomes are first
differences across the fiscal years bracketing the two snapshots.

Table~\ref{tab:esg_change} and Figure~\ref{fig:esg_change} summarize how measured ESG risk moved
within firms between 2024 and 2025. The striking pattern is in governance: measured G risk fell
for 88.9\% of firms (mean change $-1.19$), a broad-based decline consistent with a common,
methodology-driven recalibration rather than firm-specific governance improvements. Environmental
risk rose for a slight majority of firms (mean $+0.38$), social risk was roughly balanced (mean
$+0.02$), and controversy scores were unchanged for 80.1\% of firms.

\input{../scripts/R/_outputs/tab_esg_change.tex}

\begin{figure}[htbp]\centering
  \includegraphics[width=0.85\linewidth]{fig_esg_change.pdf}
  \caption{Distribution of within-firm changes in Sustainalytics E, S, and G risk,
  March 2024 to March 2025. Higher scores denote higher risk. The governance distribution is
  shifted markedly to the left (lower risk).}
  \label{fig:esg_change}
\end{figure}

% =============================================================================
\section{Empirical Strategy}\label{sec:strategy}
% =============================================================================
I work in first differences. For firm $i$ and pillar $p \in \{E, S, G\}$, let
$\Delta R_{ip} = R_{ip,2025} - R_{ip,2024}$ denote the within-firm change in measured pillar
risk, and let $\Delta Y_i$ denote the change in an outcome. The structural relationship of
interest is
\begin{equation}
  \Delta Y_i = \beta \, \Delta R_{ip} + X_i' \gamma + \varepsilon_i,
  \label{eq:structural}
\end{equation}
where $X_i$ collects predetermined firm controls (log assets, leverage, return on assets, and log
market capitalization). Ordinary least squares estimates of $\beta$ are biased if firm-specific
shocks move both measured ESG risk and outcomes---precisely the concern that motivates an
instrument.

\paragraph{A shift-share instrument from the rating overhaul.} I instrument $\Delta R_{ip}$ with
the leave-one-out industry mean of the pillar change,
\begin{equation}
  Z_{ip} = \frac{1}{|I(i)| - 1} \sum_{j \in I(i),\, j \neq i} \Delta R_{jp},
  \label{eq:instrument}
\end{equation}
where $I(i)$ is firm $i$'s industry. This is a shift-share (Bartik) instrument
\citep{GoldsmithPinkhamSorkinSwift2020_bartik, BorusyakHullJaravel2022_shift_share}: it captures the
common, industry-level component of the 2024 rating-system change---the methodology-driven shift
experienced by a firm's industry peers---purged of firm $i$'s own idiosyncratic revision.
Identification is therefore \emph{cross-industry}: I do not include industry or sector fixed effects,
because the instrument is itself industry-level variation and such effects would mechanically absorb
it. The exclusion
restriction requires that the common industry-level revision in measured risk affect firm
outcomes only through firm $i$'s own measured-risk change; industry-level confounds that move both
the average revision and outcomes are the central threat, and I return to this in
Section~\ref{sec:conclusion}.

\paragraph{Instrument strength.} Table~\ref{tab:first_stage} and Figure~\ref{fig:first_stage}
report robust first-stage $F$ statistics. The shift-share instrument is strong for governance
($F = 15.0$) and moderate for social risk ($F = 9.5$), exceeding the corresponding provisional
industry-exposure dummies, whose own first-stage $F$ is negligible for environmental and social
risk ($F = 0.03$ and $0.34$) and moderate for governance ($F = 9.8$). Environmental risk is not
strongly instrumented in these data ($F = 5.2$); I therefore report E only as an OLS association
and reserve causal language for S and G. Because the S and G first stages, while clearing the
conventional $F > 10$ rule of thumb (G) or approaching it (S), sit below the Stock--Yogo
10\%-maximal-size value of 16.4, I treat them as moderately strong and report
weak-instrument-robust inference alongside conventional standard errors
\citep{AndrewsStockSun2019_weak_iv, LeeMcCraryMoreiraPorter2022_valid_tratio}.

\input{../scripts/R/_outputs/tab_first_stage.tex}

\begin{figure}[htbp]\centering
  \includegraphics[width=0.8\linewidth]{fig_first_stage.pdf}
  \caption{Robust first-stage $F$ statistics by pillar for the shift-share instrument versus the
  provisional industry-exposure dummy. The dashed line marks the conventional $F = 10$ threshold.}
  \label{fig:first_stage}
\end{figure}

% =============================================================================
\section{Results}\label{sec:results}
% =============================================================================
Table~\ref{tab:main_results} and Figure~\ref{fig:coef_iv} report the main estimates. For each
outcome I show the two-stage least squares coefficient on the change in social and governance
risk (instrumented as in Equation~\ref{eq:instrument}) and the OLS coefficient on the change in
environmental risk. The six outcomes are controversy, leverage, asset growth, capital expenditure
intensity, net debt issuance, and share buybacks. No IV coefficient is statistically
distinguishable from zero at conventional levels. The point estimates are also economically
small relative to their outcome scales.

\input{../scripts/R/_outputs/tab_main_results.tex}

\begin{figure}[htbp]\centering
  \includegraphics[width=0.9\linewidth]{fig_coef_iv.pdf}
  \caption{Shift-share IV estimates of the effect of changes in social and governance risk on firm
  outcomes. Thick bars are 2SLS Wald 95\% confidence intervals; thin, lighter bars are
  weak-instrument-robust Anderson--Rubin 95\% intervals.}
  \label{fig:coef_iv}
\end{figure}

\paragraph{Weak-instrument-robust inference.} Because the first stages are only moderately strong,
conventional Wald intervals may understate uncertainty. Table~\ref{tab:ar_ci} reports
Anderson--Rubin (AR) confidence intervals, obtained by grid inversion of the robust AR test and
valid under weak identification \citep{AndrewsStockSun2019_weak_iv}. As expected, the
AR intervals are wider than the Wald intervals,
and noticeably so for social risk, whose first stage is weaker than governance's. In every case,
however, the AR intervals still contain zero, so the null results are not an artifact of
spuriously narrow conventional standard errors.

\input{../scripts/R/_outputs/tab_ar_ci.tex}

\paragraph{Interpretation.} Two readings are consistent with these results. The methodology-driven
revisions in measured ESG risk---most visibly the broad decline in governance risk---may simply
not be the kind of information that moves real corporate behavior within a year. Alternatively, a
single rating vintage and one year of differences may leave too little signal, given moderate
instruments, to detect modest effects. I cannot fully separate these, but the weak-instrument-robust
intervals show that the data are informative enough to rule out large effects on the outcomes
considered. To fix magnitudes, consider leverage (measured as the ratio of total debt to assets):
the governance estimate is $0.019$ (Anderson--Rubin 95\% interval $[-0.005, 0.058]$), so a one-point
increase in measured governance risk moves the leverage ratio by about $1.9$ percentage points at
the point estimate and by at most roughly $6$ percentage points at the upper bound. Scaled by the average within-firm governance-risk decline of
about $1.2$ points, the implied leverage change is on the order of $-2$ percentage points and is
bounded well below the kind of capital-structure shift that would warrant concern---small, not
merely insignificant.

\subsection*{Heterogeneity by rule exposure}
A null average effect can mask offsetting responses across firms. The SEC rule itself suggests
where to look: it phases in by filer size and binds most on climate-sensitive sectors. I therefore
interact the change in measured pillar risk with two predetermined moderators: a large-filer
indicator (above-median assets) and a climate-sector indicator (Energy, Utilities, Basic Materials,
and Industrials). I focus on the real-behavior outcomes. Because the
governance and social instruments become uninformative once interacted (the interacted first
stages are weak), I report the heterogeneity for environmental risk, where the relevant variation
is observed directly; these are ordinary-least-squares associations, not causal effects.

Table~\ref{tab:heterogeneity} and Figure~\ref{fig:heterogeneity} show that the relationship
between the change in environmental risk and real financing and investment behavior is
concentrated in exactly the rule-exposed firms. For large filers, a rise in environmental risk is
associated with higher leverage ($+0.006$, $p < 0.01$) and higher capital expenditure ($+0.005$,
$p < 0.05$); for climate-sector firms, it is associated with lower net debt issuance ($-0.005$,
$p < 0.05$). Three of these interactions survive a Benjamini--Hochberg correction at $q < 0.10$
across the interaction family, so the pattern is unlikely to be an artifact of multiple testing.
The average null therefore conceals a meaningful margin: environmental-risk changes move real
corporate behavior for the firms on which the disclosure rule bears most directly, even as they do
not on average. I interpret this as suggestive evidence on \emph{who} responds, to be sharpened with
the market outcomes and additional rating vintages discussed below.

\input{../scripts/R/_outputs/tab_heterogeneity.tex}

\begin{figure}[htbp]\centering
  \includegraphics[width=0.85\linewidth]{fig_heterogeneity.pdf}
  \caption{Interaction coefficients on the change in environmental risk times rule-exposure
  moderators (OLS), with 95\% confidence intervals. Diamonds mark interactions surviving a
  Benjamini--Hochberg correction at $q < 0.10$.}
  \label{fig:heterogeneity}
\end{figure}

% =============================================================================
\section{Conclusion}\label{sec:conclusion}
% =============================================================================
The 2024 U.S.\ climate-disclosure episode left a clear but narrow footprint on measured ESG risk:
within-firm governance risk fell for nearly nine in ten firms over the year, while environmental
and social risk and controversy scores moved little. Instrumenting the change in measured social
and governance risk with the common industry-level component of the contemporaneous Sustainalytics
rating overhaul, I find no statistically detectable pass-through to controversy or to real
financing and investment behavior in the first year, and weak-instrument-robust intervals confirm
that conclusion.

Three limitations bound the interpretation and point to the natural next steps. First,
identification rests on a single rating-methodology vintage and one year of differences, which
limits power; additional snapshots would sharpen both the first stage and the outcome estimates.
Second, the shift-share instrument identifies effects \emph{across} industries and cannot include
industry fixed effects, so any industry-level shock that moved both the average ESG-risk revision
and firm outcomes would violate the exclusion restriction---the assumption most in need of
external scrutiny. Third, the outcomes here are accounting-based; whether changes in measured ESG
risk are priced---as the green-asset-pricing literature would predict
\citep{PastorStambaughTaylor2022_green_returns, PedersenFitzgibbonsPomorski2021_esg_frontier} and as
provider rating changes appear to move returns \citep{GalemaGerritsen2025_esg_rating_changes}---is
the first-order open question, and answering it requires market data (cost of capital, liquidity, and
returns) not yet in the sample. Bringing those outcomes to bear on the same design is the immediate
priority.

% =============================================================================
\bibliographystyle{plainnat}
\bibliography{../Bibliography_base}

\end{document}
